\chapter{Conclusion}
\label{chap:conclusion}

This thesis evaluated four families of inverse-kinematics method on two pneumatic soft-robot datasets---a two-module six-chamber I-Support recording and an undocumented three-valve recording---under identical training and evaluation protocols. The methods compared were analytical piecewise-constant-curvature inverse kinematics in single- and two-segment formulations, feed-forward neural inverse kinematics in tip-only and tip-plus-mid input configurations, recurrent neural inverse kinematics in the same two input configurations, and Jacobian-learned iterative inverse kinematics. A sim-to-real transfer study using a calibrated PyElastica rod at four real-data budgets, and a four-trajectory tracking evaluation applied uniformly to every method, completed the experimental programme.

The headline quantitative findings are as follows. The analytical piecewise-constant-curvature baseline is not competitive with even a small feed-forward network on either dataset, for two distinct reasons: a constant-arc-length structural ceiling produced a $22$\,mm forward residual on Dataset 2, and a multi-valued-inverse calibration collapse produced a $51$\,mm closed-loop error on Dataset 1 despite a near-zero geometric residual. Adding the mid-marker coordinates to the Dataset 1 input reduced pressure MAE by $15\%$ on the feed-forward network and raised the coefficient of determination from $0.32$ to $0.44$; combining the mid-marker input with a recurrent architecture lifted the coefficient of determination further to $0.51$, the best Dataset 1 result reported in this thesis. The warm-started Jacobian-learned iteration reached $0.71$\,mm mean closed-loop position error on the per-sample Dataset 1 test sub-sample and $0.77$\,mm on the reference trajectories, hitting the DLS solver's $1$\,mm tolerance on every Dataset 1 trajectory target; on the under-actuated Dataset 2 the same iteration diverged on the helical spiral and lost the advantage, and a sim-to-real-fine-tuned feed-forward network was the best Dataset 2 tracker at $6.20$\,mm. The sim-to-real pre-training produced no measurable improvement in pressure-space accuracy on Dataset 2 at any real-data budget, a negative result attributable to a five-fold workspace mismatch between the I-Support-inspired simulator and the target robot rather than to the real-data sample count.

Against the published literature, the Dataset 1 sub-millimetre Jacobian result matches the order of magnitude reported by Fang et al.\ \cite{fang2022jacobian} for learned-Jacobian inverse kinematics on a different pneumatic platform, and the feed-forward-over-analytical improvement on Dataset 2 matches the factor-of-two pattern reported by Giorelli et al.\ \cite{giorelli2015nn} on a cable-driven soft arm. The sub-$5\%$ pressure-space MAPE benchmark of Ma et al.\ \cite{ma2022bp} was not reached on Dataset 2; most of the gap is attributable to the undocumented actuation units and zero-frame of Dataset 2 rather than to any limitation of the methods themselves.

Taken together, these findings support a single synthesis. The right inverse-kinematics method for a new pneumatic soft-robot arm depends on three properties of the arm and its intended use---actuation redundancy, sensing coverage, and target-domain overlap---and no single method tested here dominated every combination of those properties. Where actuation redundancy is available and partial-shape sensing is present, as on the two-module I-Support recording, the warm-started Jacobian-learned iteration drives closed-loop error to the solver tolerance; where actuation matches the output dimension and no documented digital twin is available, as on the three-valve recording, a feed-forward network with or without sim-to-real pre-training remains the best option. The three-property criterion organises the regime-specific results of the experimental programme into a compact method-selection framework for practitioners who have to choose an inverse-kinematics method for a new pneumatic soft-robot arm.

These findings operationalise the title of this thesis. \emph{AI for Control of Soft Robot by Using ML and Inverse Kinematics} has been instantiated here, per the scope fixed in Section~\ref{sec:intro_scope}, as (i) \emph{AI} and \emph{ML} in the form of four learning-based families---feed-forward regression, recurrent regression, Jacobian-learned iteration, and sim-to-real fine-tuning; (ii) \emph{Inverse Kinematics} in the form of both the analytical piecewise-constant-curvature baseline and the Jacobian-learned hybrid, which inverts an ML-trained forward map iteratively at run time and therefore combines \emph{ML} and \emph{IK} in a single architecture; and (iii) \emph{Control of Soft Robot} in its open-loop position-control sense, evaluated on the four-trajectory tracking protocol through a shared learned forward-kinematics arbiter. Within this scope, the three-property decision criterion above gives the practical shape of an answer to the title question: which AI instantiation best controls a soft robot depends on the platform's actuation redundancy, its sensing coverage, and the overlap between any simulator prior and its real operating workspace---and the experimental evidence of Chapter~\ref{chap:results} quantifies that dependence on two complementary pneumatic datasets. Closed-loop real-hardware validation of the controlled arm remains the natural next step, as discussed in Chapter~\ref{chap:future_work}.
